% =====================================================================
%  GMN3 §6 + §7.1  线算子与模空间上的全纯函数；六维 (2,0) 理论回顾
%  原文：Lamination-Paper.tex 3695--4378 行；PDF p.57--66（印刷 p.56--65）
%  编号公式 (6.1)--(6.15)、(7.1)--(7.5) 共 20 条，全部在本块内，故 \eqref 可用；
%  跨块引用（节号、式号）暂写原文编号并加 % TODO，合稿时统一改 \ref/\eqref
% =====================================================================

\section{线算子与模空间上的全纯函数}
\label{sec:linehol}

\subsection{在「达布坐标」中的展开}

我们一直在研究一个 $d=4$、$\CN=2$ 的理论，它配备一族线算子 $\CL$。现在考虑该理论住在欧几里得时空 $\IR^3\times S^1$ 上，其中圆周的半径为 $R$，并且对费米子取周期边界条件。低能下，该理论由三维 sigma 模型描述，其靶空间是超凯勒流形 $\CM_\CL$。（下文第 6.3 节% TODO \ref{sec:moreonmoduli}
将详述：对不同的 $\CL$，这些模空间彼此如何不同。）

这个 sigma 模型是 \cite{Gaiotto:2008cd} 的主要研究对象。我们先快速回顾几个基本点，以固定记号。拓扑上 $\CM_\CL$ 是 $\CB$ 上的环面纤维化。$u\in\CB$ 处的纤维，是阿贝尔规范场绕 $S^1$ 的电与磁 Wilson 线所成的环面。选定配对 $(-1)^{\inprod{\gamma,\gamma'}}$ 的一个二次加细 $\sigma:\Gamma_\CL\to\IZ_2$ 之后，该环面纤维可以等同于特征标环面 $\Hom(\Gamma_\CL,\IR/2\pi\IZ)$。我们把此环面中的元素记作 $\theta$。$\CM_\CL$ 上有复的「达布坐标」$\CX_\gamma$，它们在 \cite{Gaiotto:2008cd,Gaiotto:2009hg} 中已被详细研究。\footnote{关于超凯勒流形及其复结构的约定，以及达布坐标 $\CX_\gamma$ 的定义，我们参见 \cite{Gaiotto:2008cd}。简要总结见 \cite{Gaiotto:2009hg} 第 2 节。}
在下文中，改用坐标 $\CY_\gamma$ 会更为方便；它与 $\CX_\gamma$ 只差一个二次加细：
\begin{equation}\label{eq:Ygam-Xgam}
\CX_\gamma = \sigma(\gamma) \CY_\gamma,
\end{equation}
$\CY_\gamma$ 其实比 $\CX_\gamma$ 更典范：它们的定义不依赖于 $\sigma$ 的选取。$\CY_\gamma$ 服从一条扭乘法法则
\begin{equation}\label{eq:Y-twisted}
 \CY_\gamma \CY_{\gamma'} = (-1)^{\inprod{\gamma,\gamma'}} \CY_{\gamma+\gamma'}.
\end{equation}

现在可以让四维理论的线算子 $L_\zeta$ 绕紧化圆周 $S^1$ 而缠绕，从而定义相应的圈算子。在低能极限下，它们成为三维理论中的点算子。这样一个点算子的期望值 $\inprod{L_\zeta}$ 是模空间 $\CM_\CL$ 上的函数。如附录 B% TODO \ref{app:Holomorphy}
所说明的，$L_\zeta$ 保持 $osp(4^*\vert 2)_\zeta$ 这一事实，蕴涵 $\inprod{L_\zeta}$ 实际上是 $\CM_\CL$ 上（复结构 $J^{(\zeta)}$ 下）的\ti{全纯}函数。由此，它们可以表成 $\CY_\gamma$ 的函数。我们猜想
\begin{equation} \label{eq:trh-large-r}
\inprod{L_\zeta} = \sum_\gamma \fro(L_\zeta, \gamma) \CY_\gamma.
\end{equation}

下面我们试着为 \eqref{eq:trh-large-r}
给出论证。我们觉得下述论证是有说服力的，尽管它够不上一个证明。先从把 $\inprod{L_\zeta}$ 看作四维理论希尔伯特空间上的一个迹说起。要把这个迹显式写出，略有微妙之处：
\begin{equation}\label{eq:Lvev-trace}
\inprod{L_\zeta} = {\rm Tr}_{\CH_{u,L_\zeta}} (-1)^F e^{-2 \pi R H} e^{i \theta\cdot \CQ}\sigma(\CQ).
\end{equation}
这里有两点微妙之处。第一，$\CM_\CL$ 的环面纤维定义了一个具有非平凡单值的局部系统。$\CQ$ 是荷算子，它测量 $\Gamma$ 中的红外荷，因而取值于某个局部系统。记号 $\theta\cdot\CQ$ 是自然的求值配对，它不带单值。第二点微妙之处是：$\theta$ 同时为电 Wilson 线\emph{与}磁 Wilson 线定义了边界条件。要落实这样的边界条件，我们应当使用自对偶的形式体系。众所周知，要定义自对偶阿贝尔规范理论的路积分，就必须为相应的上同调引入一个双线性形式的二次加细——自对偶规范理论的 Dirac 量子化正建基于这个上同调之上 \cite{Witten:1996hc,Freed:2000ta,Hopkins:2002rd,Belov:2006jd}。在当下情形，我们应取 Dirac 量子化条件取值于 $H^2(M_4;\Gamma)$。当 $M_4=(\IR^3-\{\vec 0\})\times S^1$ 时，使用紧支撑上同调，我们就被引导到在迹中包含二次加细 $\sigma(\CQ)$。这就是式 \eqref{eq:Lvev-trace}
中 $\sigma(\CQ)$ 的来历。

在 $R\to\infty$ 极限下，迹 \eqref{eq:Lvev-trace}
被投影到能量最低的那些态，即带框 BPS 态上。于是可以使用理论的红外描述。对一个电荷为 $\gamma$ 的带框 BPS 态，我们预期的主要贡献形如
\begin{equation}
\sigma(\gamma) (-1)^{F} \exp(2 \pi R \, {\rm Re} (Z_\gamma / \zeta) + i \theta_\gamma).
\end{equation}
因此当 $R\to\infty$ 时
\begin{equation}\label{eq:Lzeta-asymp}
\inprod{L_\zeta} \sim \sum_\gamma \fro(L_\zeta, \gamma) \exp(2\pi R
\, {\rm Re} (Z_\gamma / \zeta) + i \tilde\theta_\gamma),
\end{equation}
其中 $e^{i \tilde\theta_{\gamma}}:= e^{i \theta_{\gamma}}\sigma(\gamma)$。这个展开与我们所想要的 \eqref{eq:trh-large-r}
很像，只不过那里写的是 $\CY_\gamma$，这里写的是 $\exp(2\pi R \, {\rm Re} (Z_\gamma / \zeta) + i \tilde\theta_\gamma)$。而这确实与 \cite{Gaiotto:2008cd} 所确定的 $\CY_\gamma$ 在 $R\to\infty$ 的渐近行为相符。那么如何延拓到有限的 $R$ 呢？我们暂且把这个展开写成
\begin{equation} \label{eq:trh-large-r-prov}
\inprod{L_\zeta} = \sum_\gamma \fro(L_\zeta, \gamma) \CY'_\gamma,
\end{equation}
其中 $\CY'_\gamma$ 是 $\CM$ 上（复结构 $\zeta$ 下）的一组全纯函数。在红外理论中，$\CY'_\gamma$ 刻画的是插入一个电荷为 $\gamma$ 的源的效果。特别地，我们预期关于紫外线算子的全部细节信息都被 $\fro(L_\zeta, \gamma)$ 吸收，从而 $\CY'_\gamma$ 不依赖于紫外线算子 $L_\zeta$ 的具体选取。如果假定线算子 $L_\zeta$ 足够多，那么我们其实可以把 \eqref{eq:trh-large-r}
看作定义 $\CY'_\gamma$ 的一个线性变换。我们想证明 $\CY_{\gamma}=\CY'_\gamma$。

回忆 \cite{Gaiotto:2008cd,Gaiotto:2009hg}：$\CX_\gamma$（从而 $\CY_\gamma$）可以被一组公理唯一刻画，这组公理表述它们的全纯性、$R\to\infty$ 与 $\zeta\to 0,\infty$ 时的渐近行为、$\zeta\to -1/\bar\zeta$ 下的实性、穿墙行为以及乘法。让我们逐条比较 $\CY'_\gamma$ 与 $\CY_\gamma$。

首先是全纯性：我们知道 $\inprod{L_\zeta}$ 在 $\CM^\zeta$ 上全纯，并且按假定在 $\zeta\in\IC^*$ 上全纯。其次，如我们已看到的，$\CY'_\gamma$ 在 $R\to\infty$ 的渐近行为就是 $\CY_\gamma$ 的渐近行为，即
\begin{equation}\label{eq:X-asymps}
\CY'_\gamma \sim \CY_{\gamma}^{\rm sf} = \exp( \pi R
 \frac{Z_\gamma}{\zeta} + i \tilde\theta_\gamma + \pi R \zeta \bar Z_{\gamma} ).
\end{equation}
对 $\zeta\to 0,\infty$ 也有类似的渐近行为，这可从超对称 Wilson 圈的形式性质结合对偶推出。用同样的策略，也可以论证 $\zeta\to -1/\bar\zeta$ 下所需的实性。

接着考虑穿墙。没有令人信服的理由预期 $\inprod{L_\zeta}$ 作为 $\CB\times\IC^*$ 上的函数会出现任何相变，因为紫外理论的真空结构没有相变，而且只要 $R>0$，我们取的始终是一个性质良好的迹类算子的迹。因此 $\inprod{L_\zeta}$ 应当\emph{不}表现出任何穿墙行为。另一方面，正如我们在第 3.4 节% TODO \ref{subsec:Halo-Wall-Crossing}
中详细看到的，带框 BPS 简并度 $\fro(L_\zeta,\gamma)$ 确实经历穿墙。现在注意：\eqref{eq:trh-large-r}
看起来极像我们在 (3.38)% TODO \eqref{eq:form-gen}
中引入的形式生成函数 $F(L)$，只是把形式变量 $X_\gamma$ 换成了 $\CM$ 上货真价实的函数 $\CY_\gamma$，并把 $y$ 取为 $y=-1$。当 $(u,\zeta)$ 越过一道墙 $\widehat{W}(\gamma_\BPS)$ 时，函数 $\CY_\gamma$ 的间断由一个辛同胚 $\CK_{\gamma_\BPS}^{\Omega(\gamma_\BPS)}$ 给出，其中 $\CK_{\gamma}$ 是 (3.36)% TODO \eqref{eq:KS-sympo}
中定义的辛同胚。因此，与从 (3.26)% TODO \eqref{eq:gen}
到 (3.34)% TODO \eqref{eq:y-min-one-gen}
的讨论相比较，我们看到：只要 $\CY'_\gamma$ 的变换方式与 $\CY_\gamma$ 完全相同，$\inprod{L_\zeta}$ 就确实不会遭受任何穿墙间断。

最后，$y=-1$ 时的海森堡关系 (3.37)% TODO \eqref{eq:Heis-Alg}
蕴涵函数 $\CY_\gamma$ 应满足扭乘法法则 $\CY_\gamma \CY_{\gamma'} =(-1)^{\langle \gamma,\gamma'\rangle} \CY_{\gamma+\gamma'}$。这与 \eqref{eq:Y-twisted}
美妙地相容。

至此我们完成了对 \eqref{eq:trh-large-r}
与 \eqref{eq:Ygam-Xgam}
的论证。论证的主要缺口在于如下断言：线算子足够多，足以把方程 \eqref{eq:trh-large-r}
反解出来。


\subsection{关于达布展开的评注}

我们想就上一小节作若干评注。

\begin{enumerate}

\item 由于指标在有质量表示上为零，人们可能会问：为什么 \eqref{eq:Lvev-trace}
不简单地是带框 BPS 态上的求和？（正如在指标定理的热核证明中、或在紧靶空间上超对称量子力学里 Witten 指标的计算中，确实会发生的那样。）
如我们已指出的，把迹中的算子作用在带框 BPS 态上，给出的是半平坦达布坐标，而不是真正的达布坐标，因而给出错误的答案。其原因与 \cite{Cecotti:1992qh} 中长篇讨论的现象非常相似：连续谱对迹作出了非零贡献。

\item 正如我们在第 7 节% TODO \ref{sec:Six-Dim-View}
中所描述的，有一类非常有趣的 $d=4$、$\CN=2$ 理论，来自 $(2,0)$ 理论在带穿刺点的黎曼面 $C$ 上的紧化。在这类情形，红外阿贝尔规范理论来自阿贝尔 $(2,0)$ 理论的自对偶 3-形式在 $\IR^4\times\Sigma$ 上的紧化，其中 $\Sigma$ 是覆盖 $C$ 的 Seiberg--Witten 曲线。在阿贝尔六维理论的路积分中，对拓扑扇区的求和由一个二次加细加权，该二次加细是 $H^3(M_6;\IZ)$ 上模二相交形式的二次加细 \cite{Witten:1996hc,Freed:2000ta,Hopkins:2002rd,Belov:2006jd}。把它应用于 $M_6=(\IR^3-\{\vec 0\})\times S^1\times\Sigma$，并使用紧支撑上同调，六维理论的这个二次加细就诱导出 $H^1_{cpt}(\Sigma;\IZ)\cong H_1(\Sigma;\IZ)$ 上模二相交形式的二次加细 $\sigma$。

\item 既然有保护自旋特征标，自然可以把 \eqref{eq:Lvev-trace}
的右端推广，引入一个与 $\CJ_3$ 共轭的变量，并定义扭迹：
\begin{equation}\label{eq:y-dep-trace}
\inprod{L_\zeta}_y := {\rm Tr}_{\CH_{u, L_\zeta}} (-1)^F e^{-2 \pi R H} (-y)^{2 \CJ_3} e^{i \theta \cdot \CQ}\sigma(\CQ)
\end{equation}
迹的被积算子仍与一个超对称 $\epsilon_{\alpha A} \CR_{\alpha}^A$ 对易。然而，\eqref{eq:y-dep-trace}
的物理解释看来与 \eqref{eq:Lvev-trace}
相当不同。原因在于：若要给 \eqref{eq:y-dep-trace}
一个三维的解释，则这个被修改过的迹要求我们在绕 $z$ 轴旋转与 R 对称性旋转之下把空间重新粘回自身。这让人想起四维背景的一种 $\Omega$ 形变，它只保持一个超荷。（近期文献可参见 \cite{Nekrasov:2010ka}，其中引有更早的文献。）
用三维场论来解释这个量并不直接，而且可能要用一个非交换的靶空间来替代 $\CM$。的确，人们可以猜测：低能下理论归结为 $\CM$ 上的一个量子力学问题，而迹投影到最低能量的态上，在那里相应的展开系数是非交换的（很像量子霍尔效应最低朗道能级中的波函数可以用非交换几何来解释）。这或许是一个富有成果的语境，可为第 3.6 节% TODO \ref{subsec:Deformed-Ring}
所讨论的线算子形变手征环给出物理解释。这是一个有意思的进一步研究方向，但本文不去追究，只在下文第 11 节% TODO \ref{sec:Quantum-Holonomy}
作一些相关的评注。

\item 最后，我们可以就如下事情作一个简单的观察：相应线算子的自旋特征标之乘积，给出了 $\CM$ 上全纯函数环的一个非交换形变，它是什么意思。离开 $y=\pm 1$ 的首个修正由泊松括号 $\{ X_\gamma,X_{\gamma'}\} = \langle \gamma,\gamma'\rangle X_\gamma X_{\gamma'}$ 给出，它与 $\CX_\gamma$ 这些函数的自然泊松括号一致（记住：超凯勒流形 $\CM$ 在每一个复结构下都是复辛的）。因此，这个非交换乘积应当与 $\CM$ 上全纯函数环的形变量子化一致。我们在第 8 节% TODO \ref{sec:Examples-of-M}
中用若干例子说明这一评注。

\end{enumerate}


\subsection{多个模空间与极大互局部格} \label{sec:moreonmoduli}

现在我们回到 $\inprod{L_\zeta}$ 所定义于其上的那些模空间的精确定义。

设 $\CL_i$ 标记简单线算子的各种不同的极大互局部组。对每个 $i$ 有一个电荷格 $\Gamma_{\CL_i}$，它们嵌在一个包含关系系统之中\footnote{这样一套包含关系在布尔代数里也被称为一个「格」（lattice）。不过，那个词即便不引起彻底的混淆，大概也会造成别扭，所以我们不用它。}
\begin{equation}
\begin{matrix}
    &  &  \Gamma_{mx} &  &   \\
    & \nearrow &   &  \nwarrow & \\
    \Gamma_{\CL_1} &  & \cdots &  & \Gamma_{\CL_n} \\
    & \nwarrow &   &  \nearrow & \\
    &  &  \Gamma  &  &   \\
\end{matrix}
\end{equation}
其中 $\Gamma_{mx}$ 是诸 $\Gamma_{\CL_i}$ 的并。注意反对称配对 $\inprod{,}$ 在 $\Gamma_{mx}$ 上不必取整数值。这尤其蕴涵：$\Gamma_{mx}$ 一般并不对应任何物理 $\CN=2$ 理论中的电荷格。

每个电荷格 $\CL$ 对应一个略有不同的物理理论。约化到三维时，这一差别变得重要：三维 sigma 模型的靶是依赖于 $\CL$ 的超凯勒流形 $\CM_\CL$。它是一个环面纤维化，$u\in\CB-\CB^{sing}$ 上方的纤维由
\begin{equation}
(\CM_{\CL})_u = {\rm Hom} ((\Gamma_{\CL})_u, \IR/\IZ).
\end{equation}
给出。我们可以类似地定义对应于格 $\Gamma$ 与 $\Gamma_{mx}$ 的空间 $\CM$ 与 $\CM_{mx}$，尽管它们并不对应任何物理的 $\CN=2$ 理论。于是有一族有限的覆盖
\begin{equation}\label{eq:covering-maps}
\begin{matrix}
    &  &  \CM_{mx} &  &   \\
    & \swarrow &   &  \searrow & \\
    \CM_{\CL_1} &  & \cdots &  & \CM_{\CL_n} \\
    & \searrow &   &  \swarrow & \\
    &  &  \CM &  &   \\
\end{matrix}
\end{equation}

事实上，这些流形彼此关系极为密切：$\CM$ 与 $\CM_{mx}$ 都是超凯勒的（正如诸 $\CM_{\CL_i}$ 一样），并且 \eqref{eq:covering-maps}
中所有覆盖映射都是局部等距。的确，它们是环面的纤维式覆盖。这一点容易从 \cite{Gaiotto:2008cd} 的构造推出：度规对纤维坐标 $\tilde\theta$ 的依赖只通过因子 $e^{i\tilde\theta_\gamma}$ 出现，其中 $\gamma\in\Gamma$ 是一个「原味」BPS 粒子的荷。因此，任何使 $e^{i\tilde\theta_\gamma}$ 对一切 $\gamma\in\Gamma$ 都不变的 $\tilde\theta$ 平移，都是一个等距。

我们相信上述图像可以延拓到奇异轨迹 $\CB^{sing}$ 之上。理由如下。在一般的点 $u^*\in\CB^{sing}$ 处，纤维化中的奇异性来自单个电荷为 $\gamma$ 的 BPS 粒子变得无质量。在这种情形，环面中有一个圆在 $u^*$ 上方缩成一点。（用 \cite{Gaiotto:2009hg} 第 13.7 节的结果，这就是由流 $\frac{d\tilde\theta_{\tilde \gamma}}{dt} = \langle \gamma, \tilde \gamma \rangle$ 所描述的那个圆。）至少当时空圆的半径 $R$ 很大时，我们可以用带荷 $\gamma$ 的粒子对朴素的大半径（半平坦）度量的单圈修正，来理解 $u^*$ 附近完整的超凯勒度量的行为。用 \cite{Gaiotto:2008cd} 的结果我们看到：若 $\gamma$ 不可约，则该圆只在纤维中的一处缩掉，即 $\exp i \tilde\theta_\gamma =1$ 处。否则若 $\gamma=Q\gamma_0$，则度量有 $Q$ 个余维二的奇点，它们——一般而言——是成族的 $A_{Q-1}$ 型奇点。用这幅纤维退化的图像，我们看到 \eqref{eq:covering-maps}
中商空间的关系得以保持。

我们将在第 7.3.1 节% TODO \ref{subsec:Hitchin-Isogenous}
再次评注这些问题，并在第 8 节% TODO \ref{sec:Examples-of-M}
中用具体例子说明这些想法。


\subsection{扭模空间的一个 vev} \label{sec:twisted-vev}

既然引入 \eqref{eq:y-dep-trace}
已经打开了潘多拉魔盒，我们应当指出：除了上面讨论的 $y=-1$ 情形之外，对一个密切相关的迹，\emph{确实}存在一种解释：
\begin{equation}\label{eq:y=1-trace}
\inprod{L_\zeta}' := {\rm Tr}_{\CH_{u,L_\zeta}} (-1)^{2I_3} e^{-2 \pi R H} e^{i \theta \cdot \CQ}
\end{equation}
其中我们选取 $u,\theta$，并用这个迹来\emph{定义}左端。同样地，四个超荷全都被保持。额外的因子 $(-1)^{2 I_3}$ 作用在非异类带框 BPS 态上是平凡的。

与 $y=-1$ 的情形一样，$\inprod{L_\zeta}'$ 可以用一个三维 sigma 模型来理解：它由我们的四维理论在 $S^1$ 上紧化得到，只是现在多了一个扭因子 $(-1)^{2 I_3}$。所得的理论和以前一样，是映入某个超凯勒流形的 sigma 模型，我们把它记作 $\tCM$。

$\tCM$ 与 $\CM$ 是什么关系？两者都可以用 \cite{Gaiotto:2008cd} 的方法来描述。那个描述中主要的非平凡成分是一组瞬子修正，它们来自四维理论的 BPS 态，其世界线缠绕紧化圆周 $S^1$。那么因子 $(-1)^{2 I_3}$ 如何改变这个故事呢？如果存在异类 BPS 态\zy{原文 exotic：指在 $su(2)_R$ 下作非平凡变换的 BPS 态，见第 3.2 节。% TODO \ref{subsec:BPS-Deg-PSC}
此处译作「异类」，以免与 singular（奇异）相混。}，则相应的瞬子修正会被这个因子影响。因此，在那种情形 $\tCM$ 会与 $\CM$ 不同。然而在没有异类 BPS 态时，我们预期 $\CM$ 与 $\tCM$ 作为超凯勒流形是同构的（尽管或许不是典范地同构）。至少在我们下面遇到的例子里，情况确实如此。

与上面同样的推理使我们能作展开
\begin{equation} \label{eq:y1}
\inprod{L_\zeta}' = \sum_\gamma \fro(L_\zeta, \gamma, y=1) \tCY_\gamma,
\end{equation}
其中 $\tCY_\gamma$ 是 $\tCM$ 上的一族函数，满足与 $\CY_\gamma$ 类似的性质。重要的一点是：与 $\CY_\gamma$ 不同，这些函数服从\ti{不带扭}的乘法法则
\begin{equation}
\tCY_{\gamma} \tCY_{\gamma'} = \tCY_{\gamma + \gamma'}.
\end{equation}

既然我们已经说过，在没有异类 BPS 态时流形 $\CM$ 与 $\tCM$ 应当重合，人们也许想知道如何看出这一等同。最显然的做法是：把 $\tCM$ 上的函数 $\tCY_\gamma$ 与 $\CM$ 上的函数 $\sigma(\gamma)\CY_\gamma$ 等同起来，其中 $\sigma$ 是某个选定的二次加细。这个等同与 $\CY_\gamma$ 及 $\tCY_\gamma$ 的间断性质相容，当且仅当 $\sigma(\gamma)=(-1)^{2 \CJ_3}$（作用在一切「原味」BPS 态上）。这一性质确实对 \cite{Gaiotto:2009hg} 中为 $A_1$ 理论所定义的典范 $\sigma$ 成立。如果强正性成立，我们就应预期更一般的理论也有类似的简单行为。


\section{六维视角}\label{sec:Six-Dim-View}

到此为止，我们的考虑只用到 $d=4$、$\CN=2$ 超对称。在本文余下部分，我们将聚焦于 \cite{Gaiotto:2009hg} 中记作 class $\CS$（取「六」之意）的那类理论。这些理论来自超共形 $ADE$ $(2,0)$ 理论在带穿刺点的黎曼面 $C$ 上的紧化。它们构成一大批丰富的例子，而六维视角将让我们能构造出一些非常有趣的线算子。对 $A_1$ 理论，我们将能够把线算子的真空期望值完全算出来。


\subsection{六维 $(2,0)$ 理论回顾}\label{subsec:6d-review}

IIB 型弦论或 M 理论中的一般性考虑，提示存在具有超共形对称性的六维定域量子场论 \cite{Witten:1995zh,Strominger:1996ac,Witten:1995em}。人们对这些理论所知甚少。关于已知结果的精彩综述，见 \cite{Seiberg:1997ax,Witten:2009at}。下面我们总结若干与本文相关的已知结果。特别地，我们聚焦于具有 $(2,0)$ 超共形对称性的理论。关于五维与六维超对称的约定之总结，见
附录 E% TODO \ref{app:Six-Super}。

张量多重态的阿贝尔理论是 $(2,0)$ 理论的例子。它们是自由场论，可以用拉氏量描述（在作出某些选择之后 \cite{Belov:2006jd}）。我们转而聚焦于相互作用理论：对它们，人们并不知道、甚至也不指望存在基本的场论自由度，当然更没有拉氏量。这些理论有 $ADE$ 分类；本节中 $\lieg$ 指某个单且单链李代数的紧实形式。

在 $A_r$ 理论的情形，可以利用 DLCQ 以及瞬子模空间上的量子力学给出一个部分定义 \cite{Aharony:1997th,Aharony:1997an}。关于 $ADE$ 理论，我们所知道的基于两条我们将当作公理来用的陈述。一旦这些理论被正确定义，它们就应当成为定理：

\begin{enumerate}

\item 当理论紧化在 $\IR^{1,4}\times S^1$ 上（$S^1$ 半径为 $R$，费米子取周期边界条件）时，长距离动力学由极大超对称的 Yang--Mills 理论支配，其规范李代数为 $\lieg$，耦合常数 $g^2_{YM}$ 正比于 $R$。\footnote{关于按我们约定给出的精确归一化，见 \cite{Gaiotto:2009hg} 的式 (3.36)。}

\item $\IR^{1,5}$ 上的理论有一个真空模空间
\begin{equation}
\CM(\IR^{1,5}) = (\IR^5\otimes \liet)/\CW,
\end{equation}
其中 $\liet$ 是 $\lieg$ 的嘉当子代数，$\CW$ 是 Weyl 群。在 $\CM(\IR^{1,5})$ 的光滑点处，低能动力学由取值于 $\liet$ 的自由 $(2,0)$ 张量多重态描述。
\end{enumerate}

当然，这两种形变必须相容。例如，五维超对称 Yang--Mills 理论的真空模空间可以等同于 $\CM(\IR^{1,5})$，因为 $so(5)_R$ 的 ${\bf 5}$ 中伴随标量的真空期望值 $\langle Y^I \rangle$ 必须可同时对角化。模空间的一个点因而被等同于 $\lieg$ 中一组有序的、可同时对角化的五元组，模掉规范不变性。正如这些 vev 在五维 SYM 中把 $\lieg$ 对称性自发破缺到嘉当子代数 $\liet$，我们应当把六维理论的「规范对称性」设想为自发破缺到取值于 $\liet$ 的自由张量多重态的阿贝尔规范群。由于这个缘故，模空间 $\CM(\IR^{1,5})$ 常被称为「库仑分支」。

此外，在 $A_r$ 理论的情形，有一个很好的直观图像：该理论是 $(r+1)$ 张重合 M5 膜的世界体，处在引力解耦极限下，并已把质心自由度除掉 \cite{Strominger:1996ac}。$so(5)$ R 对称性于是被解释为垂直于五维膜的空间中的旋转，而模空间 $(\IR^5 \otimes \IR^{r+1})/S_{r+1}$ 参数化 $r+1$ 张平行、单次缠绕的 M5 膜在横空间中的位置。单张 M5 膜的低能动力学由 $U(1)$ 张量多重态描述，标量场 $Y^I$ 表示横向涨落。模空间的一个点因而可以由标量场的真空期望值 $\langle Y^I_s \rangle$（$s=1,\dots, r+1$）指定。$A_r$ 理论是在把描述整体质心的那个张量多重态「解耦」之后得到的。

$(2,0)$ 理论可以定义在某些洛伦兹六维流形 $M_6$ 上，其配分函数可以维克转动到欧几里得六维流形。张量多重态中规范势 $B$ 的 $3$-形式场强被约束为自对偶。这蕴涵 $M_6$ 必须可定向，并且还必须配备某些额外的拓扑数据。\footnote{技术上，是第四 Wu 类的一个微分积分提升 \cite{Hopkins:2002rd}。关于这一条件的物理讨论，见 \cite{Witten:1996hc} 与 \cite{Belov:2006jd}。}
此外，为了定义费米子，它还必须配备自旋结构。

最后，我们提一下这些六维理论中一个非常微妙的方面，此处不打算深入讨论。六维理论不是标准的量子场论，而大概最好看成取值于一个可逆七维拓扑场论的六维理论。\footnote{我们在此不去细究这一概念的确切细节。感谢 D. Freed 与 E. Witten 就这一问题所作的有益讨论。}「向量值量子场论」这个词有时也被使用，因为这些理论在六维流形上的配分函数取值于一个向量空间，而不只是一个数。这一微妙之处在圆周紧化时就显现出来。我们上面说五维 SYM 的规范李代数是 $\lieg$，但要指定理论，还必须选一个李代数为 $\lieg$ 的规范群 $G$。这个规范群的选取并不唯一。必须指定额外的数据，它们决定配分函数向量空间到一维子空间的一个投影 \cite{Witten:2009at}。物理上，这样的投影或许可以由五维理论中 't Hooft 通量的一个选取来指定。这一点在 ${\cal N}=4$ SYM 的 S 对偶性质中已经可见；该理论正是六维理论的一个环面紧化：弱耦合时，在环面的「小」圆上约化会选定一个五维规范群，随后在第二个圆上约化又出现同一个规范群。S 对偶交换这两个圆，并交换对应于 Langlands 对偶规范群的两个投影。


\subsubsection{手征原初场}

$A_r$ 理论的 DLCQ 定义表明，存在生成超共形代数短多重态的定域手征原初场 \cite{Aharony:1997an,Bhattacharya:2008zy}。在圆周紧化之后，它们可以被等同于由标量场的 Casimir 不变量构造的规范不变保护算子，并在 $SO(5)$ 指标上对称且无迹。在 $A_r$ 情形，它们下落到五维 SYM 中的算子 ${\Tr} Y^{(I_1 \cdots} Y^{I_k)}$。它们的标度维数为 $2k$。看来可以合理地假定：这样的手征原初多重态不仅对 $A_r$ 存在，而且对一切 $ADE$ 理论都存在；我们在下文中作此假定。一个秩为 $d$ 的 Casimir 给出一个维数为 $2d$ 的算子，它处在 $SO(5)$ 的对称无迹表示中，带 $d$ 个指标。


\subsubsection{动力学 BPS 弦}

六维非阿贝尔理论中有动力学弦。在 $A_r$ 理论的 M 理论推导中，它们来自张开在 M5 膜之间的开 M2 膜。单个 BPS 弦态由根 $\alpha\in\Phi(\lieg)$ 标记，其中 $\Phi(\lieg)$ 是 $\lieg$ 的根系。沿弦的单位切向量为 $\hat t^M$ 的弦，在超对称代数中有中心荷 $Z_M^I = \hat t_M z^I$。弦的张力正比于 $\parallel \vec z \parallel$，其中 $\parallel \cdot \parallel$ 是 $\IR^5$ 上的欧几里得范数。库仑分支上位于点 $Y\in \IR^5 \otimes \liet$ 的真空，其弦激发有 $\vec z$ 正比于 $\langle \alpha, Y\rangle$。应当把这些弦看成类似于非阿贝尔规范理论中的「非对角」规范玻色子。

我们可以通过在 $S^1$ 上紧化并利用所预期的到五维 SYM 的约化，来核验这些断言。BPS 弦下落到五维 SYM 的 BPS 粒子。对应于单次缠绕弦的粒子，对应于荷为 $\alpha$ 的规范玻色子，质量为 $\sim R \lVert \langle \alpha, Y\rangle \rVert$。当 $R$ 小时不应有其他轻粒子（否则它们必须由五维 SYM 多重态中不存在的场来描述），所以例如不应存在由 $\lieg$ 的其他权向量标记的弦。

六维的动力学弦也在五维 SYM 中给出 BPS 弦态。在库仑分支上，可以把它们半经典地描述为垂直于弦世界面的三维空间中的磁单极子。这样的解因而由 $\Lambda_{G} = {\rm Hom}(U(1),T)$ 标记，即 $G$ 的余特征标格。对于\emph{单链}李代数，在选取 Killing 形式的归一化使根长为 $2$ 之后，$\Lambda_{cr}$ 与 $\Lambda_{r}$ 同构。由于这个缘故，五维中的弦在取规范群为 $\tilde G$（与 $\lieg$ 相伴的单连通紧群）时最容易理解。

为完整起见，我们提一下：五维 SYM 中还有由世界线横向的四个维度中的瞬子所描述的 BPS 粒子。它们的质量正比于 $n/R$，其中 $n$ 是瞬子荷，可以把它们解释为由 $S^1$ 上的紧化产生的 KK 模式。


\subsubsection{面算子}

除了手征原初场与动力学弦之外，六维 $(2,0)$ 理论及其五维 SYM 紧化中还有面算子。本节的主要目标正是用它们来定义线算子，而后者是本文余下部分的研究重点。

我们先给出面算子的一个一般定义，它与线算子的定义密切类似：$\IR^{1,1}\subset\IR^{1,5}$ 上的面算子，其横空间是 $\IR^4$。度规共形等价于 $AdS_3\times S^3$，其等距群的李代数由 $so(2,2)\oplus so(4)$ 给出。我们把嵌入的 $\IR^{1,1}$ 上的面算子定义为：在 $AdS_3\times S^3$ 上强加一个共形不变的边界条件所得的结果。注意 $so(2,2)\oplus so(4) \cong (sl(2,\IR) \oplus su(2)) \oplus (sl(2,\IR) \oplus su(2))\cong so(4^*) \oplus so(4^*)$。

与线算子一样，我们将聚焦于那些保持若干 R 对称性与若干超对称的特殊面算子。为说明其动机，先考虑库仑分支上低能张量多重态理论中的面算子。若 $\Sigma$ 是 $M_6$ 中的定向曲面，则可以定义自对偶场的和乐\footnote{在下文的讨论中，为简单起见我们假定 $B$ 在拓扑上平凡。对一般的时空与拓扑非平凡的 $B$ 场，我们应改用 Cheeger--Simons 微分特征标。}
\begin{equation}\label{eq:self-dual-holon}
h_v(\Sigma) = \exp[ 2\pi i \int_{\Sigma} (v, B) ]
\end{equation}
其中 $v \in \Lambda_{wt}$ 属于 $\tilde G$ 的特征标群，即权重格。我们强加这一限制，是因为在拓扑非平凡的流形上，$B$ 场的大规范变换取 $B \to B + \omega$ 的形式，其中 $\omega$ 是整体定义的闭 2-形式，其周期取值于 $\Lambda_{cr}$。注意这里以及第 7 节% TODO \ref{sec:Six-Dim-View}
余下部分中，$\Sigma$ 与红外 Seiberg--Witten 曲线\emph{毫无}关系！

在某些情形下，这个和乐可以提升为一个超对称面算子。为简单起见，考虑规范群为 $U(1)$ 的单个张量多重态。于是考虑
\begin{equation}\label{eq:single-tm-susy}
\exp\left[ 2\pi i \int_{\Sigma} B + \kappa n^I Y^I \vol(\Sigma)\right]
\end{equation}
其中 $\vol(\Sigma)$ 是由诱导度规给出的 $\Sigma$ 上的体积形式，$\kappa$ 是常数，并且不失一般性可以假定 $\vec n$ 是欧几里得 $\IR^5$ 中的单位向量，其分量为 $n^I$。设 $\xi^\alpha$（$\alpha = 1,2$）是 $\Sigma$ 上的局部坐标系。那么超对称 $\epsilon_i^r Q_r^i$ 湮灭此算子，当且仅当
\begin{equation}\label{eq:surface-susy-cond}
\epsilon_i^r \left( \frac{d\xi^\alpha \wedge d \xi^\beta \p_\alpha X^M \p_\beta X^N}{\vol(\Sigma) }(\gamma_{MN})_r^{~~s} \delta^i_{~~j} + \kappa (n^I \Gamma^I)^i_{~~j} \delta_{r}^{~~s} \right) =0
\end{equation}
其中 $X^M(\xi)$ 表示曲面到 $M_6$ 中的嵌入。为保持超对称，此式必须对\emph{常数}的未破缺超对称 $\epsilon_i^r Q_r^i$ 成立。对平坦曲面与常数 $n^I$，当 $\kappa=\pm 1$ 时有一半超对称被保持。更一般地（类似于超对称 Yang--Mills），$\Sigma$ 可以是弯曲曲面，$n^I$ 也可以变动。一个在下文会变得重要的例子是曲面为 $\IR \times \wp$ 的情形，其中 $\wp$ 是比如 $12$ 平面中的一条曲线。若把 R 对称性空间分解为 $\IR^5 = \IR^2 \oplus \IR^3$，把 $\IR^2$ 加项等同于 $12$ 平面，并取 $n^I$ 为 $\wp$ 的单位切向量，则保持四分之一的超对称。一般而言，R 对称性被方向 $n^I$ 破缺到 $so(4)$。在比如 $01$ 方向的平面的特殊情形，我们可以用与线算子的 $osp(4^*\vert 2)_\zeta$ 讨论密切类似的方式描述所保持的超对称。我们通过沿曲面所在平面作反射、再作垂直于 $n^I$ 的平面内的 R 对称性反射，来定义 $osp(8^*\vert 4)$ 的一个对合。此对合的不动点定义一个超共形代数，我们可以把它记作 $(osp(4^*\vert 2) \oplus osp(4^*\vert 2))_{\vec n}$。

现在可以把 \eqref{eq:self-dual-holon}
与 \eqref{eq:single-tm-susy}
结合起来，在阿贝尔张量多重态理论中定义超对称面算子 $\IS_{\vec n}(v,\Sigma)$。于是很自然地猜想：非阿贝尔理论中也存在相应的面算子 $\IS_{\vec n}(\CR,\Sigma)$，它与 $\lieg$ 的表示 $\CR$ 相伴，并保持超共形对称性 $(osp(4^*\vert 2) \oplus osp(4^*\vert 2))_{\vec n}$。像往常一样，我们可以用两个可用的形变来核验这一断言。当 $Y$ 的大期望值把规范对称性强力破缺时，这些面算子将被半经典地描述为阿贝尔面算子 $\IS_{\vec n}(v,\Sigma)$ 的求和，其中 $v$ 遍历表示 $\CR$ 的权重。作为第二项核验，我们考虑在半径为 $R$ 的圆周上紧化后出现的与五维 SYM 的关系。若 $\Sigma = S^1\times \wp$，其中 $\wp\subset M_5$ 是一条路径，则非阿贝尔面算子的红外极限应当是
\begin{equation}\label{eq:5d-line-op}
L_{\vec n}(\CR, \wp) = {\Tr}_{\CR} P\exp \left[ \int_\wp  A + \kappa n\cdot R Y ds \right]
\end{equation}
这些确实就是五维 SYM 中的 BPS 线算子。最后，非阿贝尔面算子的存在与动力学 BPS 弦的存在密切相连。例如，可以从一个 $su(r+1)$ 理论出发，走到部分库仑分支，在那里理论约化为 $su(r) \times u(1)$ 理论。将存在带 $u(1)$ 阿贝尔因子之荷的动力学弦，其张力当标量期望值趋于无穷时趋于无穷，而 $u(1)$ 因子解耦。在此极限下，动力学弦成为面算子 $\IS_{\vec n}(v,\Sigma)$，其中 $\vec n$ 是被送到无穷远的 $Y$ 的方向。

众所周知，把诸如 \eqref{eq:5d-line-op}
这样的线算子插入路积分，可以解释为把背景理论与一个 $0+1$ 维量子力学系统（由余伴随轨道的量子化定义）耦合起来 \cite{Alekseev:1988vx}。类似地，面算子可以由一个定义在 $\Sigma$ 上的 $1+1$ 维量子场论与背景理论耦合来表示。在这一解释中要注意：正如人们从 $AdS_3/CFT_2$ 对应中所熟悉的，等距的李代数 $so(2,2)\cong sl(2,\IR)_l \oplus sl(2,\IR)_r$ 可以解释为分别作用在左行与右行自由度上的共形对称代数。这一解释可以延拓，使我们能把紫外理论中未破缺的超共形对称性解释为 $osp(4^*\vert 2)_l \oplus osp(4^*\vert 2)_r$。此外，如我们已看到的，在红外，由 $\IS_{\vec n}(\CR,\Sigma)$ 下落的阿贝尔面算子由表示 $\CR$ 的权重 $v$ 标记。这提示：定义在 $\Sigma$ 上的那个量子场论应当有以 $v$ 标记的真空。

动力学弦可以终止在面算子上，它们在 $\Sigma$ 上的二维量子场论中定义畴壁。由根 $\alpha \in \Phi(\lieg)$ 标记的弦，若 $v_1 - v_2 = \pm \alpha$，就是真空 $v_1$ 与 $v_2$ 之间的畴壁。这样的构型可以是 $\frac{1}{4}$ BPS 的，其方式与动力学弦之间的汇点成为 $\frac{1}{4}$ BPS 构型的方式十分相似。一个典型的例子是静态构型：所有弦与面算子都位于一个平面上，中心荷向量 $\langle \alpha_i,Y\rangle$ 或 $\vec n$ 也都位于中心荷空间中的一个平面上，并且两个平面中的斜率相同。注意，若关心的是 $\frac{1}{4}$ BPS 构型，则面算子可以允许比平面中的直线更一般的形状，办法是让 $\vec n$ 沿算子变化，跟随切向量的变化。
